%-*-latex-*-
%-*-latex-*-
\newcommand\COURSE{ciss240}
\newcommand\ASSESSMENT{q1006}
\newcommand\ASSESSMENTTYPE{Quiz}
\newcommand\POINTS{\textwhite{xxx/xxx}}

\input{myquizpreamble}
\input{yliow}
\input{\COURSE}
%-*-latex-*-

%-*-latex-*-

%-*-latex-*-
\renewcommand\TITLE{\ASSESSMENTTYPE \ \ASSESSMENT}

\renewcommand\EMAIL{}
\renewcommand\AUTHOR{eagraham1@cougars.ccis.edu}

\textwidth=6in
\begin{document}
\topmatterthree

This is a closed-book, no compiler, 5 minute quiz.

%------------------------------------------------------------------------------
\nextq
The following is the graph of $y = x^2 \sin(x^3) + 10$:
\begin{python}
from latextool_basic import *
p = FunctionPlot(domain='0:3', height='1.5in')
p.add('x * sin(x**3) + 10', color='red', python=1, pin='below', num_points=1000,
pin_message='')
print(p)
\end{python}
Using exactly 1000 rectangles, compute the area under the above curve
from $x = 0.5$ to $x = 2.5$.
The area of each rectangle must be computed using the method in the notes
(i.e., use the left endpoint of the rectangle to determine the area of the rectangle).
The printout of the area must be in fixed point format with 5 decimal
places.
Skeleton code is provided.
\\
\ANSWER
\begin{answercode}
#include <iostream>
#include <iomanip>
#include <cmath>

int main()
{
    double area = 0.0;                // the required area
    double a = 0.5, b = 2.5;          // bounds
    int n = 1000;                     // number of rectangles
    double dx = (b - a) / n;          // width of each rectangle

    for (int i = 0; i < n; ++i)       
    {
        double x = a + i * dx;        
        double y = x * x * sin(x * x * x) + 10; 
        area += y * dx;
    }
    std::cout << std::fixed << std::setprecision(5) << area << std::endl;
    return 0;
}
\end{answercode}

%------------------------------------------------------------------------------
\nextq
The following is the graph of $y = x^2 \sin(x^3)$:
\begin{python}
from latextool_basic import *
p = FunctionPlot(domain='0:3', height='1.5in')
p.add('x**2 * sin(x**3)', color='red', python=1, pin='below', num_points=1000,
pin_message='')
p.add('0', color='blue', python=1, pin='below', num_points=10,
pin_message='')
print(p)
\end{python}
Notice that there's a root at approximately $x = 1.7$.
Use exactly 1000 equally spaced points from $a = 1.6$ to $b = 2.0$ to
compute an approximation to this root.
The printout of the value $x$ of the root
must be in fixed point format with 5 decimal places.
Skeleton code is provided.
\\
\ANSWER
\begin{answercode}
#include <iostream>
#include <iomanip>
#include <cmath>

int main()
{
    double area = 0.0;                // the required area
    double min_y;                     // the running minimum y for given x 
    double a = 0.5, b = 2.5;          // bounds
    int n = 1000;                     // number of points
    double dx = (b - a) / n;          // width of each rectangle

     for (int i = 0; i <= n; ++i)
    {
        double x = a + i * dx;        
        double y = x * x * sin(x * x * x); 

        if ((y > 0 && y < min_y) || (y < 0 && -y < min_y)) 
        {
            min_y = (y > 0) ? y : -y;
            root_x = x;
        }
    }

    std::cout << std::fixed << std::setprecision(5) << area << std::endl;
    return 0;
}

\end{answercode}

%------------------------------------------------------------------------------
%-*-latex-*-
\newpage
\input{instructions}
\end{document}

